\section{Background}

\subsection{From Content-Level Risks to Action-Level Hazards}

Conventional LLM systems are primarily exposed to \emph{content-level} failure
modes, including misleading generation, privacy-sensitive text leakage, and
overreaching recommendations. In these systems, harmful output is usually
bounded by language itself: the model ``says'' something wrong, unsafe, or
inappropriate. In contrast, Computer-Use Agents (CUAs) possess
environment-execution capabilities, allowing model outputs to be translated
into concrete operations across software interfaces. As a result, classical
LLM risks are not merely preserved but operationally amplified into
\emph{action-level} hazards with higher real-world impact.

\begin{itemize}
\item \textbf{Expanded input attack surface.} CUA perception is no longer limited
to a single prompt channel. Web pages, pop-up dialogs, emails, local
documents, chat messages, and desktop notifications can all function as
instruction carriers. This multi-channel exposure increases susceptibility to
indirect prompt injection, where malicious or irrelevant interface content
steers the agent away from user intent while appearing contextually
plausible [3--4, 16--18]. Unlike direct jailbreaking, these attacks are
embedded in ordinary task environments, making detection harder and
raising the likelihood of silent objective drift.

\item \textbf{Privilege and toolchain connectivity.} CUAs are often integrated
with persistent browser sessions, local file systems, terminal/script
execution, and authenticated account states. This creates a continuous path
from ``influence'' to ``execution'': an attacker-controlled instruction can
progress from content manipulation to privileged operations, such as data
exfiltration, unauthorized transaction steps, or unsafe command invocation
[4, 13--15, 33--36]. Therefore, the security boundary is no longer the model
response surface; it extends to every connected tool and credential context.

\item \textbf{Amplification in long-horizon workflows.} CUA tasks are frequently
multi-step, cross-application, and stateful. As horizon length grows, small
reasoning errors, permission mistakes, or policy violations can compound and
be automatically propagated across downstream actions. This temporal
amplification makes incidents more difficult to audit, attribute, and roll
back, especially when intermediate decisions are implicit or weakly logged
[26--31, 60--61]. Consequently, reliability and governance for CUAs must be
evaluated as process integrity over time, not only as single-turn response
quality.
\end{itemize}

Taken together, these properties indicate a categorical shift in risk: CUAs
should be modeled as socio-technical execution systems rather than text-only
assistants. This motivates security controls that combine prompt-level
defenses with runtime policy enforcement, least-privilege tool access,
action confirmation boundaries, and end-to-end traceability.

\subsection{Ecosystem-Specific Security: MCP and Skills}

Beyond generic CUA risks, two ecosystem components deserve explicit treatment:
\emph{MCP-based tool integration} and \emph{skill-based capability reuse}.
Both improve extensibility, but both also introduce additional trust
dependencies and compositional failure modes.

\textbf{MCP trust boundary risks.}
MCP-style integration allows an agent to access external tools/resources
through server endpoints. If endpoint identity, transport integrity, or tool
schema constraints are weak, a malicious (or compromised) MCP server can
inject manipulated outputs that appear operationally valid. This can convert
context poisoning into privileged side effects, including unauthorized file
access, leakage of sensitive workspace context, and hidden execution of unsafe
operations via abstract tool calls. In practice, MCP security is therefore not
only a connectivity issue; it is a cross-boundary execution integrity issue.

\textbf{Skill supply-chain and execution risks.}
Skills package prompts, scripts, and reusable workflows. This packaging model
reduces engineering overhead, but it creates plugin-like supply-chain risk.
Untrusted or poorly reviewed skills may contain over-privileged steps,
dangerous command templates, covert data egress logic, or prompt constructions
that degrade policy compliance. Even benign skills can become unsafe through
version drift and undocumented behavioral changes. Accordingly, skill adoption
requires provenance validation, version pinning, permission scoping, and
runtime policy checks rather than one-time functional testing.

\textbf{Compositional risk in long-horizon execution.}
The highest-impact failures often occur when skills orchestrate multi-step
actions across several MCP servers. Local checks may pass at each hop while
the end-to-end chain still violates user intent or security policy. This
visibility gap complicates accountability because different components may log
different abstractions of the same action. Hence, CUA assurance must evaluate
the entire execution graph (inputs, tool calls, credentials, state changes,
and side effects) under adversarial conditions, not isolated components.

\textbf{Governance implications.}
These properties motivate concrete controls in future CUA governance
architecture: authenticated MCP channels, strict tool I/O typing and policy
enforcement, least-privilege credential delegation, sandboxed skill execution,
high-risk action confirmation, and tamper-evident audit trails for incident
reconstruction. In short, MCP servers and skills should be treated as
first-class security principals, not neutral implementation details.
 
